Dựa trên những bằng chứng về sự tồn tại của cấu trúc phụ thuộc phi tuyến và hiện tượng hội tụ rủi ro cực đoan ở trên, nghiên cứu tiến hành khớp dữ liệu thực tế vào ba mô hình Copula tiêu biểu: Gaussian Copula (đại diện cho phân phối chuẩn, không có phụ thuộc đuôi), t-Student Copula (đại diện cho phụ thuộc đuôi đối xứng) và Clayton Copula (đại diện cho phụ thuộc đuôi dưới bất đối xứng).\\
Để xác định mô hình phản ánh sát nhất với thực tế thị trường, nghiên cứu sử dụng tiêu chuẩn khoảng cách Cramér-von Mises. Cụ thể, thuật toán tiến hành tính toán tổng bình phương sai số (MSE) giữa hàm Copula thực nghiệm và các hàm Copula lý thuyết được mô phỏng với kích thước mẫu 10.000 kịch bản. Mô hình có giá trị MSE nhỏ nhất sẽ được lựa chọn.

\begin{table}[H]
    \centering
    \caption{Kết quả kiểm định độ phù hợp dựa trên trung bình bình phương sai số (MSE)}
    \vspace{0.2cm}
    \begin{tabular}{|l|c|c|}
        \hline
        \textbf{Mô hình Copula} & \textbf{Trung bình bình phương sai số (MSE)} & \textbf{Xếp hạng} \\
        \hline
        Gaussian Copula & 0.0542 & 2 \\ % Gắn số Gaussian vào đây
        t-Student Copula (df=4) & 0.0369 & \textbf{1} \\ % Gắn số t-Student vào đây
        Clayton Copula & 0.1412 & 3 \\ % Gắn số Clayton vào đây
        \hline
    \end{tabular}
    \label{tab:gof_copula}
\end{table}
\begin{flushleft}
\textbf{Nhận xét từ Bảng \ref{tab:gof_copula}:}
\begin{itemize}
    \item Kết quả đo lường chỉ ra rằng t-Student Copula có mức sai số thấp nhất (0.0369), vượt trội hơn hẳn so với Gaussian Copula (0.0542) và Clayton Copula (0.1412). Kết quả này hoàn toàn nhất quán với những phát hiện từ phân tích tương quan đuôi trước đó: cấu trúc phụ thuộc của các cổ phiếu công nghệ không tuân theo phân phối chuẩn, mà có đặc tính đuôi dày và mức độ đồng biến cao trong các trạng thái thị trường cực đoan.
    \item Việc mô hình Gaussian Copula có sai số lớn và cho thấy sự hạn chế đáng kể trong quá trình mô hình hóa, đồng thời khẳng định giới hạn của các phương pháp đo lường tuyến tính truyền thống. Trong khi đó, Clayton Copula (vốn chỉ nhạy bén với cấu trúc đuôi dưới) cũng không đạt được độ khớp tối ưu, cho thấy tính chất phụ thuộc của nhóm tài sản này mang xu hướng đối xứng hơn ở cả hai vùng cực đoan.
    \item Từ cơ sở thực nghiệm vững chắc này, t-Student Copula được lựa chọn làm mô hình nền tảng. Hàm Copula này sẽ đóng vai trò là lõi cấu trúc phụ thuộc cho thuật toán mô phỏng Monte Carlo ở bước tiếp theo, phản ánh sát với đặc tính biến động thực tế của thị trường tài chính.
\end{itemize}
\end{flushleft}